%% Chapter 4, Section 4.3: Implementation of MCMC

\section{Implementation of Markov Chain Monte Carlo Methods}
\label{sec:4.3}

\subsection{Initialization Bias and Burn-In}
\label{sec:4.3.1}

MCMC chains are usually initialized outside stationarity (if we could sample the initial draw
from the target we would not be doing MCMC!). To reduce the initialization bias caused by the
effect of the starting value, the first $M$ draws can be discarded. Estimation is then based on the
states visited after time $M$:
\[
  \calI_f[h] \approx \frac{1}{N - M} \sum_{n=M+1}^{N} h(X^{(n)}).
\]
The initial phase up to time $M$ is called the transient phase or burn-in period. How do we decide
on the length of the burn-in period? A first step would be to examine the output of the chain by
eye. This is a very crude method, but is very quick and cheap. However, this method should be
followed up by a more sophisticated analysis.

\begin{figure}[htbp]
\centering
\includegraphics[width=0.75\textwidth]{figures/fig_04_2}
\caption{\textit{The MCMC chain is typically started outside stationarity, and it may take several
iterations to reach a region of high target density.}}
\label{fig:4.2}
\end{figure}

\subsection{Assessing Convergence}
\label{sec:4.3.2}

How to assess convergence has been a fiercely discussed topic in the literature. Some of the
disagreement is due to the fact that there are different types of convergence. First, there is the
issue of whether the distribution of the chain is close to the stationary distribution. Second, there
is the issue of how well the chain explores the state space. Generally, chains that do not explore
the state space well (we say, mix slowly) tend to approach stationarity more slowly. But note
that even if the chain is at stationarity, ergodic averages of a slow mixing chain typically result
in inaccurate estimates, a third issue that is discussed under the general keyword of convergence.
All three issues are related to each other but are not exactly the same, which can be confusing.
For example, a chain with a bi-modal stationary distribution may not have reached equilibrium
yet, as it only ever has visited the region around one of the modes. Now suppose we are interested
in estimating a function that is equal for both modes. Then, it can happen that the corresponding
ergodic average converges much faster than the chain itself.

There are two general approaches to assess convergence to stationarity: convergence rate
calculation and convergence diagnostics.
\begin{enumerate}
\item \textbf{Convergence rate calculation}: When we looked at the speed of convergence for the independence sampler, we essentially performed a convergence rate calculation. The advantage
of this approach is, of course, that it is exact. In general, these calculations can be difficult
and will only apply to specific cases. Because they are exact methods, they can be overly
pessimistic as they take into account any worst-case scenario even if it has a very small
probability of occurring. At times they can be so pessimistic that the suggested burn-in
period simply becomes impractical.

\item \textbf{Convergence diagnostics}: This is the most widely used method. Convergence diagnostics
examine the output of the chain and try to detect a feature that may suggest divergence
from stationarity. They are usually relatively easy to implement. However, while they may
indicate that convergence has not been reached, they are not able to guarantee convergence.
\end{enumerate}
We will now examine some of the most commonly used diagnostics.

\medskip
\noindent\textbf{Geweke's Method.}
Geweke proposed a diagnostic method based on spectral analysis of time
series. Consider two subsequences of a run of your MCMC, one consisting of states immediately
after burn-in and the other consisting of states at the very end of the run. Suppose that the run
is of length $N$ and we want to examine if a burn-in of $M$ steps is sufficient. Consider ergodic
averages
\[
  \calI_f^A[h] := \frac{1}{N_A} \sum_{n=M+1}^{M+N_A} h(X^{(n)}),
  \qquad
  \calI_f^B[h] := \frac{1}{N_B} \sum_{n=N-N_B}^{N} h(X^{(n)})
\]
with $M + N_A < N - N_B$. If the chain has converged, then the two ergodic averages should be
similar. How similar they should be can be quantified using an appropriate central limit theorem.
The variance of ergodic averages under stationarity is determined by the spectral density of the
time series $\{h(X^{(n)})\}_{n=1}^\infty$, defined as
\[
  S(\omega) = \frac{1}{2\pi} \sum_{k=-\infty}^{\infty} \gamma_k \exp(ik\omega),
\]
where $\gamma_k = \Cov\!\bigl(h(X^{(0)}), h(X^{(k)})\bigr)$ is the $k$-lag autocovariance of the time series (see Definition
A.35 in Appendix A). The following asymptotic result holds. If the ratios $N_A/N$ and $N_B/N$ are
fixed with $(N_A + N_B)/(N - M) < 1$, then, under stationarity and as $N \to \infty$,
\[
  \frac{\calI_f^A[h] - \calI_f^B[h]}{\sqrt{\frac{1}{N_A}\widehat{S}_A(0) + \frac{1}{N_B}\widehat{S}_B(0)}}
  \to \Normal(0, 1),
\]
where $\widehat{S}_A(0)$ and $\widehat{S}_B(0)$ are spectral estimates of the variances. We can use the above asymptotics to test whether the means of the two sequences are equal subject to variation. Geweke
originally suggested taking the values $N_A = N/10$ and $N_B = N/2$.

\medskip
\noindent\textbf{Gelman and Rubin's Method.}
This method uses several chains started in initial states that are
over-dispersed compared to the stationary distribution. The idea is that if there is initialization
bias, then the chains will be close to different modes while under the stationary regime the chains
should behave similarly. If the chain tends to get stuck in a mode, then the path appears to be
experiencing stable fluctuations, but the chain has not converged yet. In this case, we speak of
\textit{meta-stability}, which is often caused by multi-modality of the stationary distribution. Gelman
and Rubin apply concepts from ANOVA techniques (and thus rely on Normal asymptotics). We
assume each of a total of $J$ chains are run for $2N$ iterations where the first $N$ iterations are
classified as burn-in. We first compute the variance of the means of the chains:
\[
  B = \frac{1}{J-1} \sum_{j=1}^{J} \Bigl(\calI_f^j[h] - \bar{\calI}_f^j[h]\Bigr)^2.
\]
Here we run $J$ chains, $\calI_f^j[h]$ is the ergodic average of the $j$-th chain based on the last $N$ iterations,
and $\bar{\calI}_f^j[h]$ is the average of the ergodic averages of the $J$ chains. This is interpreted as the between
chain variance. As $N \to \infty$ the between chain variance will tend to zero.

Then, we compute the within chain variance $W$, that is, the mean of the variance of each
chain:
\[
  W = \frac{1}{J} \sum_{j=1}^{J} \frac{1}{N-1} \sum_{n=N+1}^{2N}
      \Bigl(h(X_j^{(n)}) - \calI_f^j[h]\Bigr)^2,
\]
where $X_j^{(n)}$ denotes the $n$-th iterate of the $j$-th chain. We calculate the weighted average of both
variance estimates
\[
  V = \frac{N-1}{N}\, W + B
\]
to obtain an estimate of the target variance. We then monitor $R = \sqrt{V/W}$, which is called the
potential scale reduction factor. $R$ tends to be larger than one and will converge towards one as
stationarity is reached. Using Normal asymptotics, one can show that $R$ has an $F$-distribution
and the hypothesis of $R = 1$ can be tested. As a rule of thumb, if the 0.975-quantile is less than
1.2, no lack of convergence has been detected.

\medskip
\noindent\textbf{Raftery and Lewis' Method.}
In this diagnostic, a posterior quantile $p$ and an acceptable
tolerance $r$ for error in computing $p$ are specified. In addition, the user specifies the desired
probability $\alpha$ of obtaining an estimate for $p$ within the prescribed error tolerance. Raftery and
Lewis suggested a way to estimate the number $N$ of iterations and the burn-in period $M$ that are
necessary to satisfy the prescribed conditions.
